%!TEX TS-program = xelatex
% Bu dosya sjuoz_word2latex.py tarafından otomatik üretildi.
\documentclass[9pt]{extarticle}
\usepackage{multicol}
\raggedcolumns  % sütun boyu doldurulmaz: başlık ve paragraf aralarında büyük boşluk açılmaz
\usepackage[top=0.5in, bottom=0.75in, right=0.75in, left=0.75in, headheight=0.25in, includehead]{geometry}
\usepackage{fontspec}
\usepackage{amsmath}
\usepackage{amssymb}
\usepackage{graphicx}
% Denklem sütuna sığmazsa: önce 2, sonra 3 satıra bölünmüş biçimi; o da sığmazsa küçültülür
\newsavebox{\sjuozeqbox}
\newlength{\sjuozeqmax}
\newcommand{\sjuozeqset}{\setlength{\sjuozeqmax}{\linewidth}\addtolength{\sjuozeqmax}{-2.5em}}
\newcommand{\sjuozeqput}{\ifdim\wd\sjuozeqbox>\sjuozeqmax
  \typeout{SJUOZ-EQ-SCALED: \theequation; \the\sjuozeqmax; \the\wd\sjuozeqbox}% kontrol raporu için
  \resizebox{\sjuozeqmax}{!}{\usebox{\sjuozeqbox}}\else\usebox{\sjuozeqbox}\fi}
\newcommand{\sjuozfit}[3]{\sjuozeqset\sbox{\sjuozeqbox}{$\displaystyle #1$}%
  \ifdim\wd\sjuozeqbox>\sjuozeqmax\sbox{\sjuozeqbox}{$\displaystyle #2$}\fi
  \ifdim\wd\sjuozeqbox>\sjuozeqmax\sbox{\sjuozeqbox}{$\displaystyle #3$}\fi\sjuozeqput}
\newcommand{\sjuozfitone}[1]{\sjuozeqset\sbox{\sjuozeqbox}{$\displaystyle #1$}\sjuozeqput}
\usepackage{xcolor}
\usepackage{array}
\usepackage{tabularx}
\usepackage{xltabular}
\usepackage{pdflscape}
\usepackage{booktabs}
\usepackage{multirow}
\usepackage{caption}
\usepackage{subcaption}
\usepackage{titlesec}
\usepackage{fancyhdr}
\usepackage{qrcode}
\usepackage{xurl}
\usepackage{lastpage}
\usepackage{indentfirst}
\usepackage{eso-pic}
\usepackage{tikz}
\usepackage[colorlinks=true, urlcolor=blue, linkcolor=black, citecolor=black]{hyperref}

\titleformat{\section}{\large\bfseries\centering}{\thesection.}{1.0em}{}
\titleformat{\subsection}{\large\bfseries}{\thesubsection.}{0.5em}{}
\titleformat{\subsubsection}{\large\bfseries}{\thesubsubsection.}{0.5em}{}
\titlespacing*{\section}{0pt}{8pt}{6pt}
\titlespacing*{\subsection}{0pt}{8pt}{4pt}
\titlespacing*{\subsubsection}{0pt}{6pt}{3pt}

\newcommand{\titlesize}{\fontsize{12}{10}}
\renewcommand{\normalsize}{\fontsize{9pt}{9pt}\selectfont}
\renewcommand{\large}{\fontsize{10pt}{10pt}\selectfont}
\renewcommand{\small}{\fontsize{8pt}{8pt}\selectfont}
\DeclareCaptionLabelFormat{sjuozparen}{#1 (#2)}
\captionsetup[table]{labelfont=bf, justification=raggedright, singlelinecheck=false}
\captionsetup[figure]{labelfont=bf, justification=centering, singlelinecheck=false}
\captionsetup[sub]{labelformat=parens, labelfont=bf, font=small, justification=centering}
\definecolor{sjuozblue}{HTML}{C1D1EB}
\definecolor{colorabs}{HTML}{E2EAF4}

% Yazı tipleri: Times New Roman / Garamond yoksa benzerleri kullanılır
\IfFontExistsTF{Times New Roman}{\setmainfont{Times New Roman}}{%
  \IfFontExistsTF{TeX Gyre Termes}{\setmainfont{TeX Gyre Termes}}{%
    \IfFontExistsTF{Liberation Serif}{\setmainfont{Liberation Serif}}{}}}
\IfFontExistsTF{Garamond}{\newfontfamily{\garamond}{Garamond}}{%
  \IfFontExistsTF{EB Garamond}{\newfontfamily{\garamond}{EB Garamond}}{\newcommand{\garamond}{\rmfamily}}}

\newcommand{\sjuozdoi}{https://doi.org/10.25271/sjuoz.2026.14.4.2005}
% İlk sayfadaki bilgi kutusunun sayfa altından yüksekliği ve metin için ayrılan boşluk
\newcommand{\firstfooterbottom}{0.4in}
\newcommand{\firstpagereserve}{1.4in}
% Sadece ilk sayfada metin alanını kısalt; ilk sayfa basıldıktan sonra eski haline döndür
\makeatletter
\newlength{\sjuozfulltextheight}
\AtBeginDocument{%
  \setlength{\sjuozfulltextheight}{\textheight}%
  \global\advance\textheight by -\firstpagereserve
  \global\@colht\textheight \global\@colroom\textheight \global\vsize\textheight
  \AddToHookNext{shipout/after}{\global\textheight\sjuozfulltextheight}}
\makeatother

% Başlık ve altlık (header/footer)
\pagestyle{fancy}
\fancyhf{}
\renewcommand{\headrulewidth}{0pt}
\lhead{\large \textit{Khalid \& Yousif}}
\rhead{\small{\garamond{\textbf{SJUOZ|VOL14|OCT 2026|Pxxx-xxx}}}}
\lfoot{\raisebox{8pt}{\textit{DOI:}\href{\sjuozdoi}{\textit{\small{\sjuozdoi}}}}}
\cfoot{\raisebox{8pt}{\thepage}}

\fancypagestyle{firstpage}{%
  \fancyhf{}
  \renewcommand{\headrulewidth}{0pt}
  \rhead{\IfFileExists{SJUOZ.png}{\includegraphics[width=1.04in]{SJUOZ.png}}{\textbf{SJUOZ}}}
  \lhead{\colorbox{sjuozblue}{%
      \begin{minipage}[b]{0.83\textwidth}
        \vspace{0.15cm}
        \hspace{-3pt}\ \normalsize \noindent \garamond{Science Journal of University of Zakho}\\
        \normalsize \textit{\garamond{DOI: \sjuozdoi}}
        \vspace{0.2cm}
      \end{minipage}}}
}

% İlk sayfanın altındaki bilgi kutusu (sayfaya mutlak konumla yerleştirilir)
\newcommand{\sjuozfirstfooter}{%
  \begin{minipage}[b]{\textwidth}
      \rule{\textwidth}{1.5pt} \\[0.15cm]
      \begin{minipage}[t]{0.2\textwidth}\hspace{1pt} \small Access this article online \\
        \vspace{0.1cm}
        {\hypersetup{urlcolor=black}\qrcode[height=2.6162cm]{https://sjuoz.uoz.edu.krd/index.php/sjuoz/article/view/2005}}
      \end{minipage}\begin{minipage}[t]{0.4\textwidth}
        \small DOI:\href{\sjuozdoi}{\sjuozdoi} \\ \\[0.5cm]
        Printed ISSN 2663-628X; \\
        Electronic ISSN 2663-6298
      \end{minipage}\begin{minipage}[t]{0.4\textwidth}\small \textbf{Science Journal of University of Zakho} \\
        \textbf{Vol. 14, No. 04, pp. xxx-xxx, October-2026} \\[0.5cm]
        This is an open access article under a CC BY-NC-SA 4.0 license \\
        (\href{https://creativecommons.org/licenses/by-nc-sa/4.0/}{https://creativecommons.org/licenses/by-nc-sa/4.0/})
      \end{minipage} \\[0.15cm]
      \rule{\textwidth}{1.5pt} \\[0.15cm]
      \centering \normalsize \thepage
  \end{minipage}}

% ORCID simgesi: orcid.png varsa o, yoksa çizilmiş yeşil "iD" simgesi
\definecolor{orcidgreen}{HTML}{A6CE39}
\newcommand{\orcidicon}{\IfFileExists{orcid.png}{\includegraphics[width=9pt]{orcid.png}}{%
  \tikz[baseline=-0.7ex]{\fill[orcidgreen] (0,0) circle (4.5pt);
  \node[text=white, font=\fontsize{4.5pt}{4.5pt}\selectfont\bfseries, inner sep=0pt] at (0,0) {iD};}}}
\newcommand{\orcidlink}[1]{\href{#1}{\orcidicon}}

\begin{document}
\setcounter{page}{1}
\thispagestyle{firstpage}
\AddToShipoutPictureFG*{\put(\LenToUnit{0.75in},\LenToUnit{\firstfooterbottom}){\sjuozfirstfooter}}

\noindent \large \textbf{Original Article} \\[-5pt]
\begin{center} \titlesize \textbf{SUBFIGURES PLACED IN DIFFERENT WAYS} \end{center}
\begin{center} \large \textbf{Shilan Khalid \textsuperscript{1,*} and Bawar Yousif \textsuperscript{1}}
\end{center}
\noindent
\large \textsuperscript{1}Department of Computer Science, University of Zakho, Duhok, Iraq.\\

\begin{center} \textsuperscript{*}\textit{Corresponding author, E-mail:}\href{mailto:shilan.khalid@uoz.edu.krd}{shilan.khalid@uoz.edu.krd}
\end{center}
\rule{\textwidth}{1.5pt}
\begin{minipage}[t]{0.175\textwidth}
  \hspace{0cm} \vspace{0cm}
\end{minipage} \begin{minipage}[t]{0.8\textwidth}
\large \textbf{ABSTRACT}\end{minipage} \\[-6pt]
\rule{\textwidth}{1.5pt} \\
\begin{minipage}[t]{0.175\textwidth}
  \vspace{1cm} Received \\ xxx \\[0.5cm] Accepted \\ xxx \\[0.5cm] Published \\ xxx \\[0.25cm]
\end{minipage} \colorbox{colorabs}{\begin{minipage}[t]{0.81\textwidth}
  \large{Authors place sub-figures in Word in several ways: inside a borderless table with labels under the pictures, as several pictures in one paragraph, or as pictures in separate paragraphs.
  \\[2pt]
  \normalsize \textbf{Keywords:} subfigures; layout.
  }
\end{minipage}}
\\[-10pt]

\normalsize
\noindent \rule{\textwidth}{0.75pt}
\begin{multicols}{2}
\section{INTRODUCTION}
\hyperlink{fig:1}{Figure 1} uses a layout table, \hyperlink{fig:2}{Figure 2} has two pictures in one paragraph and \hyperlink{fig:3}{Figure 3} has two pictures in separate paragraphs.

\noindent\begin{minipage}{\linewidth}
\centering\hypertarget{fig:1}{}
\captionsetup{type=figure}
\begin{subfigure}[t]{0.47\linewidth}\centering
\includegraphics[width=\linewidth, height=0.3\textheight, keepaspectratio]{figures/Figure1a.IMG}
\caption{Original image}
\end{subfigure}%
\hfill
\begin{subfigure}[t]{0.47\linewidth}\centering
\includegraphics[width=\linewidth, height=0.3\textheight, keepaspectratio]{figures/Figure1b.IMG}
\caption{Detected mouth}
\end{subfigure}%
\par\vspace{2pt}
\caption{Pictures placed with a layout table.}
\end{minipage}
\par\vspace{6pt}

\noindent\begin{minipage}{\linewidth}
\centering\hypertarget{fig:2}{}
\captionsetup{type=figure}
\begin{subfigure}[t]{0.47\linewidth}\centering
\includegraphics[width=\linewidth, height=0.3\textheight, keepaspectratio]{figures/Figure2a.IMG}
\caption{}
\end{subfigure}%
\hfill
\begin{subfigure}[t]{0.47\linewidth}\centering
\includegraphics[width=\linewidth, height=0.3\textheight, keepaspectratio]{figures/Figure2b.IMG}
\caption{}
\end{subfigure}%
\par\vspace{2pt}
\caption{(a) Left picture. (b) Right picture.}
\end{minipage}
\par\vspace{6pt}

\noindent\begin{minipage}{\linewidth}
\centering\hypertarget{fig:3}{}
\captionsetup{type=figure}
\begin{minipage}[t]{0.47\linewidth}\centering
\includegraphics[width=\linewidth, height=0.3\textheight, keepaspectratio]{figures/Figure3a.IMG}
\end{minipage}%
\hfill
\begin{minipage}[t]{0.47\linewidth}\centering
\includegraphics[width=\linewidth, height=0.3\textheight, keepaspectratio]{figures/Figure3b.IMG}
\end{minipage}%
\par\vspace{2pt}
\caption{Two pictures in separate paragraphs.}
\end{minipage}
\par\vspace{6pt}

\section{CONCLUSIONS}
All pictures must stay with their captions (\hyperlink{ref:1}{Khalid, 2025}).

\section*{REFERENCES}
\begin{list}{}{
    \setlength{\leftmargin}{1.5em}
    \setlength{\itemindent}{-1.5em}
    \setlength{\parsep}{0pt}
    \setlength{\itemsep}{6pt}
}
\item \hypertarget{ref:1}{}Khalid, S. (2025). Placing pictures in Word. \textit{Journal of Layout}, 3(1), 1–5.

\end{list}
\end{multicols}

\end{document}
